\section{Introduction}\label{sec:intro}
\enlargethispage{2pt}
Principal component analysis (PCA) estimates the eigensystem of a covariance matrix.  Its geometric and probabilistic foundations go back to \citet{Pearson1901} and \citet{Hotelling1933}; standard accounts and a modern review are given by \citet{Jolliffe2002} and \citet{JolliffeCadima2016}.  In the classical regime, the dimension is fixed and the sample size tends to infinity.

When dimension and sample size grow proportionally, the Mar\v{c}enko--Pastur law and spiked-covariance theory describe bulk spectra, extreme eigenvalues, and eigenvector phase transitions \citep{MarchenkoPastur1967,Johnstone2001,BaikEtAl2005,BaikSilverstein2006,Paul2007}.  The present paper instead emphasizes the high-dimension, low-sample-size (HDLSS) regime $d\gg n$, where the sample covariance has rank at most $n-1$.  Its nonzero eigenvalues are more naturally studied through the $n\times n$ dual Gram matrix than through the $d\times d$ covariance matrix; see \citet{HallMarronNeeman2005}, the extension of the geometric representation by \citet{AhnEtAl2007}, \citet{JungMarron2009}, and the surveys of \citet{AoshimaYata2017,AoshimaEtAl2018}.

\citet{YataAoshima2012} developed the noise-reduction methodology (NRM), which subtracts an estimate of the high-dimensional noise floor from each empirical eigenvalue.  The cross-data-matrix methodology (CDM) of \citet{YataAoshima2010} instead splits the sample and crosses independent halves, replacing the positive within-sample noise shift by a mean-zero cross-noise term.  Related HDLSS developments include non-Gaussian PCA consistency \citep{YataAoshima2009}, intrinsic-dimension estimation \citep{YataAoshima2010Intrinsic}, first-component inference \citep{IshiiEtAl2016}, and empirical-eigenstructure bias analysis \citep{WangFan2017}.  The power-spiked analysis of \citet{YataAoshima2013} places these procedures in a common asymptotic framework.  Sparsity offers a complementary route to direction recovery, as illustrated by \citet{JohnstoneLu2009}; here we study dense directions and explicit noise-floor corrections.

The phrase ``under a sub-Gaussian assumption'' is not by itself precise enough for a quadratic-form proof.  A tail bound for every linear functional does not automatically provide Hanson--Wright concentration for $x^\top A x$.  We consequently distinguish directional sub-Gaussianity from an independent-factor representation and use the latter in the main results.

The central point is informational.  Increasing $d$ creates many bulk coordinates, so their aggregate Gram contribution can become nearly deterministic.  It does not create additional observations of the low-dimensional signal scores.  When $n$ is fixed, the score covariance remains random; even perfect bulk removal need not consistently estimate a population spike.

\subsection{Contributions}
The paper makes five contributions.
\begin{enumerate}
\item It gives a nonasymptotic decomposition of the dual covariance.  The bulk error is expressed by $\tau_0=\tr(\Lambda_0)$, $q_0=\norm{\Lambda_0}_{\F}$, and $b_0=\norm{\Lambda_0}_{\op}$.
\item It identifies the exact finite-score target of sequential NRM.  The target for component $i$ contains $-(r-i)^{-1}\sum_{j=i+1}^m s_j$, which need not vanish at fixed $n$.
\item Conditional on the correct signal rank $m$, it analyzes block NRM (bNRM), which estimates the common bulk floor from only the lower empirical eigenvalues.
\item It derives a comparable high-probability CDM bound and separates the cancellation of the mean shift from the variability introduced by sample splitting.
\item It reports deterministic, reproducible numerical experiments that distinguish population error from finite-score-oracle error rather than conflating the two targets.
\end{enumerate}

The results exchange distributional generality for explicit probabilities.  The independent-factor sub-Gaussian assumption is stronger than the fourth-moment conditions used in parts of the original NRM/CDM literature.  The claims below concern theorem form, proof transparency, and interpretation; they do not assert that every correction displayed here is absent from all previous covariance-denoising work.
